% ============================================================
%  preamble.tex  —  IBO Biology Platform
%  Minimal academic style. pdfLaTeX compatible.
% ============================================================

% --- Page geometry ---
\usepackage[a4paper,
            top=2.5cm, bottom=2.8cm,
            left=3cm, right=2.5cm]{geometry}

% --- Typography ---
\usepackage[T1]{fontenc}
\usepackage[utf8]{inputenc}
\usepackage{lmodern}
\usepackage{microtype}
\usepackage{setspace}
\setstretch{1.12}

% --- Mathematics ---
\usepackage{amsmath}
\usepackage{amssymb}
\usepackage{amsfonts}
\usepackage{mathtools}
\usepackage{bm}
\usepackage{siunitx}
\sisetup{per-mode=symbol}

% --- Chemistry ---
\usepackage[version=4]{mhchem}

% --- Graphics ---
\usepackage{graphicx}
\usepackage{xcolor}
\usepackage{tikz}
\usetikzlibrary{calc}

% --- Floats and captions ---
\usepackage{float}
\usepackage{caption}
\captionsetup{font=small, labelfont=bf, labelsep=period, width=0.92\textwidth}

% --- Tables ---
\usepackage{booktabs}
\usepackage{tabularx}
\usepackage{array}

% --- Lists ---
\usepackage{enumitem}
\setlist[enumerate,1]{label=\textbf{(\alph*)}, itemsep=2pt}
\setlist[enumerate,2]{label=\roman*.,  itemsep=1pt}
\setlist[itemize]{itemsep=2pt}

% --- Framed environments (for solutions, notes) ---
\usepackage{mdframed}

% --- Colour palette: one accent, otherwise near-black ---
\definecolor{accent}{HTML}{1A2E4A}    % dark navy
\definecolor{probback}{HTML}{F5F6F7}  % near-white box background
\definecolor{stargold}{HTML}{A07820}  % muted gold
\definecolor{rulecol}{HTML}{AAAAAA}   % light rule colour
\definecolor{codebg}{HTML}{F8F8F8}
\definecolor{codeframe}{HTML}{CCCCCC}
\definecolor{codekw}{HTML}{00397A}
\definecolor{codestr}{HTML}{B83200}
\definecolor{codecomment}{HTML}{4A7A4A}
\definecolor{codenumber}{HTML}{999999}

% --- Hyperlinks ---
\usepackage{hyperref}
\hypersetup{
    colorlinks = true,
    linkcolor  = accent,
    citecolor  = accent,
    urlcolor   = accent,
    pdftitle   = {IBO Biology Platform},
}

% --- Code listings ---
\usepackage{listings}
\usepackage{lstautogobble}

\lstdefinestyle{python}{
    language          = Python,
    backgroundcolor   = \color{codebg},
    basicstyle        = \ttfamily\small,
    keywordstyle      = \color{codekw}\bfseries,
    stringstyle       = \color{codestr},
    commentstyle      = \color{codecomment}\itshape,
    numberstyle       = \tiny\color{codenumber},
    numbers           = left,
    stepnumber        = 1,
    numbersep         = 8pt,
    frame             = single,
    framesep          = 4pt,
    rulecolor         = \color{codeframe},
    breaklines        = true,
    showstringspaces  = false,
    tabsize           = 4,
    captionpos        = b,
    autogobble        = true,
    morekeywords      = {as, with, None, True, False,
                         np, plt, odeint, linspace, exp, log, pi, array},
}
\lstset{style=python}

% --- tcolorbox (problems only) ---
\usepackage{tcolorbox}
\tcbuselibrary{breakable, skins}

% Problem box: clean, numbered, minimal
\newtcolorbox[auto counter, number within=section]{problem}[1]{%
    breakable,
    enhanced,
    colback    = probback,
    colframe   = accent,
    boxrule    = 0.5pt,
    arc        = 1pt,
    fonttitle  = \bfseries\sffamily\small,
    coltitle   = white,
    colbacktitle = accent,
    title      = {Problem~\thetcbcounter.\enspace #1},
    attach boxed title to top left={yshift=-1.5mm, xshift=3mm},
    boxed title style={arc=1pt, boxrule=0pt},
    top        = 5pt, bottom = 5pt,
    left       = 6pt, right  = 6pt,
}

% Solution: simple left-rule indent, no background
\newenvironment{solution}{%
    \begin{mdframed}[
        linewidth    = 1.5pt,
        linecolor    = accent!40,
        topline      = false,
        rightline    = false,
        bottomline   = false,
        leftmargin   = 4pt,
        innerleftmargin = 10pt,
        innertopmargin  = 4pt,
        innerbottommargin = 4pt,
        backgroundcolor = white,
    ]%
    \noindent\textit{Solution.}\enspace\ignorespaces
}{%
    \end{mdframed}%
}

% Historical note: small italic, thin left rule
\newenvironment{historicalnote}[1]{%
    \begin{mdframed}[
        linewidth    = 2pt,
        linecolor    = rulecol,
        topline      = false,
        rightline    = false,
        bottomline   = false,
        leftmargin   = 4pt,
        innerleftmargin = 10pt,
        innertopmargin  = 3pt,
        innerbottommargin = 3pt,
        backgroundcolor = white,
    ]%
    {\small\sffamily\bfseries Historical note: #1}\\[3pt]%
    \small\itshape\ignorespaces
}{%
    \end{mdframed}%
}

% Definition: bold label, thin left rule
\newenvironment{definition}[1]{%
    \begin{mdframed}[
        linewidth    = 1.5pt,
        linecolor    = accent!60,
        topline      = false,
        rightline    = false,
        bottomline   = false,
        leftmargin   = 4pt,
        innerleftmargin = 10pt,
        innertopmargin  = 3pt,
        innerbottommargin = 3pt,
        backgroundcolor = white,
    ]%
    \noindent\textbf{Definition (#1).}\enspace\ignorespaces
}{%
    \end{mdframed}%
}

% Key remark: just bold inline label, no box
\newenvironment{keyinsight}{%
    \par\medskip\noindent\textbf{Remark.}\enspace\ignorespaces
}{\par\medskip}

% Model limitations: just bold inline label
\newenvironment{limitation}{%
    \par\medskip\noindent\textbf{Limitations of this model.}\enspace\ignorespaces
}{\par\medskip}

% Summary: simple ruled block
\newenvironment{summary}{%
    \begin{mdframed}[
        linewidth = 0.5pt, linecolor = accent!50,
        leftmargin = 0pt, innerleftmargin = 8pt,
        innertopmargin = 6pt, innerbottommargin = 6pt,
        backgroundcolor = probback,
    ]%
    \noindent\textbf{Summary.}\enspace\ignorespaces
}{%
    \end{mdframed}%
}

% Code lab: just a listing caption prefix
\newenvironment{codelab}[1]{%
    \par\medskip\noindent\textbf{Code.}\enspace\textit{#1}\par\smallskip
}{\par}

% ============================================================
%  Difficulty stars — plain \ifnum, no loops
% ============================================================
\usepackage{pifont}

\newcommand{\starfilled}{{\color{stargold}\ding{72}}}
\newcommand{\starempty} {{\color{stargold!40}\ding{73}}}

\newcommand{\stars}[1]{%
  \ifnum#1>0\starfilled\else\starempty\fi
  \ifnum#1>1\starfilled\else\starempty\fi
  \ifnum#1>2\starfilled\else\starempty\fi
  \ifnum#1>3\starfilled\else\starempty\fi
  \ifnum#1>4\starfilled\else\starempty\fi
}

% ============================================================
%  Discipline tags — plain \fcolorbox, safe everywhere
% ============================================================
\newcommand{\disciplinetag}[1]{%
  {\setlength{\fboxsep}{2pt}\setlength{\fboxrule}{0.4pt}%
   \fcolorbox{accent!50}{accent!8}{%
     \textcolor{accent}{\sffamily\tiny\bfseries\vphantom{Ag}#1}}}%
}

% ============================================================
%  Section / chapter headings — clean, minimal
% ============================================================
\usepackage{titlesec}

\titleformat{\chapter}[block]
  {\normalfont\LARGE\bfseries\sffamily\color{accent}}
  {Module~\thechapter.}{0.6em}{}[\vspace{2pt}\rule{\textwidth}{0.4pt}]

\titleformat{\section}
  {\normalfont\large\bfseries\sffamily\color{accent}}
  {\thesection}{0.7em}{}

\titleformat{\subsection}
  {\normalfont\normalsize\bfseries\sffamily}
  {\thesubsection}{0.6em}{}

\titleformat{\subsubsection}
  {\normalfont\normalsize\itshape}
  {}{0em}{}

% ============================================================
%  Header / footer — minimal
% ============================================================
\usepackage{fancyhdr}
\pagestyle{fancy}
\fancyhf{}
\fancyhead[RO,LE]{\small\sffamily\thepage}
\fancyhead[LO]{\small\sffamily\nouppercase{\rightmark}}
\fancyhead[RE]{\small\sffamily\nouppercase{\leftmark}}
\renewcommand{\headrulewidth}{0.3pt}
\renewcommand{\footrulewidth}{0pt}

% ============================================================
%  Maths macros
% ============================================================
\newcommand{\kB}{k_{\mathrm{B}}}
\newcommand{\NA}{N_{\mathrm{A}}}
\newcommand{\dG}{\Delta G}
\newcommand{\ddG}{\Delta\Delta G}
\newcommand{\dH}{\Delta H}
\newcommand{\dS}{\Delta S}
\newcommand{\Kd}{K_{\mathrm{d}}}
\newcommand{\Ka}{K_{\mathrm{a}}}
\newcommand{\Keq}{K_{\mathrm{eq}}}
\newcommand{\vmax}{V_{\mathrm{max}}}
\newcommand{\Km}{K_{\mathrm{m}}}
\newcommand{\conc}[1]{[\mathrm{#1}]}
\newcommand{\od}[2]{\frac{\mathrm{d}#1}{\mathrm{d}#2}}
\newcommand{\pd}[2]{\frac{\partial #1}{\partial #2}}
\newcommand{\e}{\mathrm{e}}
\newcommand{\abs}[1]{\left|#1\right|}

% ============================================================
%  Text abbreviations
% ============================================================
\newcommand{\Ecoli}{\textit{E.~coli}}
\newcommand{\invitro}{\textit{in vitro}}
\newcommand{\invivo}{\textit{in vivo}}
\newcommand{\etal}{\textit{et al.}}

% End of preamble.tex
